\section{Data}
\label{sec:data}

\paragraph{The chain}
We scan every block from block \ScanStartHeight{}, before Lightning's first channels, to block \CutoffHeight{}, the last
block mined before 1~November~2026 00:00~UTC, with our own Bitcoin Core node. For every block we keep its fee
statistics (\texttt{getblockstats}: the size-weighted percentiles of fee rates, total fees and total weight). A full
index of the chain built from the same node, with every transaction's locktime and every input's witness, gives the
HTLC expiries and the funding fees used in \S\ref{sec:htlc} and \S\ref{sec:lifecycle}.
On the blocks from 2021, which we also queried from the node directly, its fee totals agree with
\texttt{getblockstats} on every block.

\paragraph{Public channels}
Our list of public channels joins our own gossip since May~2023, mempool.space's record of closed channels and their
close reasons~\cite{mempoolspace}, and the lnresearch gossip archive for the earlier
years~\cite{lnresearchtopology}. A short channel ID encodes the block, transaction and output of the channel's funding
output, so the list resolves to \NFundingOutpoints{} funding outpoints, each to the transaction that spends it: the
close. The list holds few of the channels that closed in 2018: the gossip archive and mempool.space's record become
dense only in January~2019, and only about 1\% of 2018's force closes are of listed channels. For 2018 we therefore
report force closes only, listed or not, which the commitment fingerprint finds completely; public, cooperative,
force-share and per-channel series start in 2019, and January~2019 is still incomplete.

\paragraph{Closes and their outputs}
A close is cooperative or forced; a close that funds a new channel in the same transaction is a splice and is left
out. Force closes of channels that are not in the list are found by the fingerprint of the commitment transaction
alone (\S\ref{sec:methods}); we call them \emph{unlisted}. Up to the cutoff the data hold \NPublicCloses{} closes of
public channels (\NPublicMutual{} cooperative, \NPublicForce{} forced) and \NUnlistedForce{} unlisted force closes.
For each of the \NOutputs{} outputs of these \NForceAll{} force closes we record the transaction that spent it, if one
had by the cutoff (\NOutputsSpent{} did), and classify the spend from its witness script: the broadcaster's delayed
sweep, the peer's sweep, an HTLC transaction or claim, an anchor spend by its owner or by a third party, or a penalty.
Statistics about what happens to a close's outputs use only closes confirmed at least 4,032 blocks (four weeks) before
the cutoff, so that their outputs had time to be spent.

\paragraph{Validation}
We check the closes two ways. First, the chain index above is a separate extraction of the chain, by a different
program: in it, the transaction that spends each public channel's funding output is the close we recorded, for all
\ValLakeCompared{} closes that could be compared, and its shape, with or without the commitment fingerprint, matches
our force or cooperative label in every case. Second, mempool.space gives a close reason for
\ValCoverageYearMinPct--\ValCoverageYearMaxPct\% of the public channels that closed in each year from 2019 to 2025,
and none for 2026. It labels closes from the chain as well, so agreement with it is a check
of consistency, not independent validation. Where it gives a reason, for \ValCompared{} of our public closes
(\ValCoveragePct\%), our label agrees on \ValAgreePct\%.

\paragraph{Client labels}
For the appendix we attach to each public close the implementations at its two ends, using the labels and
classifier of our first paper on the day of the close~\cite{kumar2026clientdiversity}. They exist from May~2023 and
inherit that paper's accuracy.

\paragraph{What the chain does not show}
The chain does not show which party broadcast a force close, so every cost is per channel, both parties together. It
shows fee bumps that confirmed, not those that were tried and failed. Private channels appear only when they force
close: their cooperative closes look like ordinary payments, so their number and their force-close rate are unknown.
\S\ref{sec:limitations} lists the consequences.
